\section*{Chapter \ref{ch:qm:probability-at-speed} --- Probability at Speed}

\begin{solution}{exo:qm:probability-at-speed:1}
$M_1 = 13 + 4 \times (364 - 13)/51 = 13 + 27.53 = 40.53$, against 35 before the deal.
\end{solution}

\begin{solution}{exo:qm:probability-at-speed:2}
(a) Yes: whether it has happened is known at each instant. (b) No: the low is known only at
the close. (c) Yes: a deterministic time is a \omterm{def:qm:probability-at-speed:stopping}{stopping time}. (d) Yes, provided the forecast is
itself adapted to the observed information.
\end{solution}

\begin{solution}{exo:qm:probability-at-speed:3}
$b/(a+b) = 3/4$. The stopped walk is bounded, so optional stopping gives $\E[S_\tau] = 0$:
$\tfrac34 \times 1 - \tfrac14 \times 3 = 0$. Winning often is paid for by losing more when
it loses.
\end{solution}

\begin{solution}{exo:qm:probability-at-speed:4}
The revisions are orthogonal with equal variances, so the last 30 minutes resolve $30/390 =
7.7\%$ of the variance.
\end{solution}

\begin{solution}{exo:qm:probability-at-speed:5}
$\E[Z] = e^{-c^2/2}\E[e^{cX}] = 1$ and $Z > 0$, so $\mathbb Q$ is a probability measure
equivalent to $\P$. Its \omterm{def:qm:probability-at-speed:charfn}{characteristic function} is $e^{\iu uc - u^2/2}$
(\cref{ex:qm:probability-at-speed:shift}): $X \sim \mathcal N(c,1)$ under $\mathbb Q$, so
$\E^{\mathbb Q}[X] = c$ and $\mathbb Q(X > c) = \tfrac12$.
\end{solution}

\begin{solution}{exo:qm:probability-at-speed:6}
$\E_4[F_5] = \E_4[M_5] = M_4$, so $\E[F_5 - F_4 \mid \mathcal F_4] = (1 - a)(M_4 - 35) = \frac{1 -
a}{a}(F_4 - 35)$: the revision is predictable, and the slope is $(1 - a)/a = 0.667$ at $a = 0.6$.
The same slope appears on the previous revision $F_4 - F_3 = a(M_4 - M_3)$, because $M_4 - 35$ is a
sum of orthogonal revisions of which $M_4 - M_3$ is one.
\end{solution}

\begin{solution}{exo:qm:probability-at-speed:7}
0.661 with a $t$-statistic of 45.9, against 0.667 in theory; the honest forecaster's slope is
$-0.005$.
\end{solution}

\begin{solution}{exo:qm:probability-at-speed:8}
A mean of zero is necessary, not sufficient: the under-reacting card forecaster also has
revisions averaging zero. A \omterm{def:qm:probability-at-speed:martingale}{martingale}'s revisions must be unpredictable from everything
known before them, starting with its own past revisions, and the test is a regression of
revisions on earlier information, with standard errors that allow for the clustering of
volatility (chapter~11) and a correction for the many signals one tries (chapter~12).
\end{solution}

\begin{solution}{pb:qm:probability-at-speed:1}
\textbf{1.} $\Var(C) = 1.80^2 + 0.30^2 = 3.33$; standard deviation USD~1.82.
\textbf{2.} Before the news, $M_t = S_t$, the price observed at $t$ (the remaining increments
and $J$ have mean zero); after it, $M_t = S_t + J$.
\textbf{3.} It is the \omterm{def:qm:probability-at-speed:doob}{Doob martingale} $\E_t[C]$ of an integrable variable.
\textbf{4.} $J$.
\textbf{5.} $\P(|J| > 0.14) = 2(1 - \Phi(0.14/0.30)) = 64.1\%$.
\textbf{6.} The revisions are orthogonal (\cref{prop:qm:probability-at-speed:orthogonal}), so
their variances add, and $M_0$ is a constant.
\textbf{7.} $150/390 \times 3.24/3.33 = 37.4\%$.
\textbf{8.} $0.09/3.33 = 2.7\%$.
\textbf{9.} $(0.09 + 10/390 \times 3.24)/3.33 = 5.2\%$: 2.7\% at the news and 2.5\% in the last
ten minutes.
\textbf{10.} 5.23\% in the 20\,000 simulated days.
\textbf{11.} No: its revision over the last ten minutes contains $J/2$, which its 15:50
revision $J/2$ predicts.
\textbf{12.} The two revisions are $J/2$ and $J/2 + $ ten minutes of increments; the
correlation is $0.15^2/(0.15\sqrt{0.15^2 + 0.0831}) = 0.46$.
\textbf{13.} Slope $\Cov/\Var = 0.15^2/0.15^2 = 1$; the simulation gives 0.98 with a
$t$-statistic of 72 over 20\,000 days.
\textbf{14.} With $t = r\sqrt{n - 1}/\sqrt{1 - r^2}$, $t = 2$ needs $n = 4(1 - r^2)/r^2 + 1 =
16$ days.
\textbf{15.} The trade earns $C - (S + J/2)$ times the sign of $J$, whose mean is $\E|J|/2 =
0.5 \times 0.30\sqrt{2/\pi} = 0.120$: 12 cents a share.
\textbf{16.} The last ten minutes then carry $3.24 \times 30/410 = 0.237$, and the share is
$(0.09 + 0.237)/3.33 = 9.8\%$.
\textbf{17.} No, if both standard deviations scale with the price: the shares are ratios of
variances.
\textbf{18.} What is resolved after the last moment one can act is risk one carries, not
information one can use; a trader deciding at 15:50 faces 5.2\% of the day's variance,
not 100\%.
\textbf{19.} Named result: \emph{the fair value of the close}: in the model, 5.2\% of the
variance of the closing price is resolved from the 15:50 news to the close (2.7\% by the news
itself), and 9.8\% if the last ten minutes are three times as volatile.
\textbf{20.} It promises that the revisions cannot be predicted from the forecaster's own
information, so nobody with that information can trade profitably against it; it promises
nothing about accuracy.
\end{solution}

\begin{solution}{iq:qm:probability-at-speed:1}
With unlimited time and credit you stop at $+1$ almost surely, but the \omterm{def:qm:probability-at-speed:stopping}{stopping time} is
unbounded and the losses before it are not uniformly integrable; with any finite credit $b$
you win $+1$ with probability $b/(b+1)$ and lose $b$ otherwise: expected value zero.
\iqlookfor{optional stopping and why its conditions matter.}
\end{solution}

\begin{solution}{iq:qm:probability-at-speed:2}
An adapted integrable process whose \omterm{def:qm:probability-at-speed:condexp}{conditional expectation} of any future value is its
current value. A stock price has a positive expected return under the real-world measure,
so it is a \omterm{def:qm:probability-at-speed:martingale}{submartingale} there; discounted by the money-market account it is a \omterm{def:qm:probability-at-speed:martingale}{martingale}
under the risk-neutral measure (chapter~5).
\iqlookfor{the definition and the role of the measure.}
\end{solution}

\begin{solution}{iq:qm:probability-at-speed:3}
For $j < k$, $\E[D_jD_k] = \E[D_j\E_{k-1}[D_k]] = 0$. A forecast that is an honest \omterm{def:qm:probability-at-speed:condexp}{conditional
expectation} has revisions no one can predict from its past revisions; if a regression of
today's revision on yesterday's has a slope, the forecast is inefficient.
\iqlookfor{taking out what is known, and the testable consequence.}
\end{solution}

\begin{solution}{iq:qm:probability-at-speed:4}
The competitor's second revision is predictable from its first, so its forecast is not a
\omterm{def:qm:probability-at-speed:martingale}{martingale}: it under-reacts. Trade against its stale price at 15:50 in the direction of the
news; in the chapter's model this earns half the mean absolute news, 12 cents a share.
\iqlookfor{predictable revisions mean an exploitable forecast.}
\end{solution}

\begin{solution}{iq:qm:probability-at-speed:5}
$\E[X \mid Y]$ is the best predictor among all functions of $Y$; the linear regression is the
best among affine ones. They coincide when $\E[X \mid Y]$ is affine, as it is when $(X, Y)$ is
jointly Gaussian.
\iqlookfor{projections onto nested subspaces; the Gaussian case.}
\end{solution}

\begin{solution}{iq:qm:probability-at-speed:6}
Store every published value with its timestamp; each day regress revisions over fixed
horizons on the previous revision and on other information available at the earlier time
(order-book imbalance, recent trades), with robust standard errors, and report the variance
shares by time of day. It lies if horizons overlap (serially correlated errors), if the
volatility is heteroskedastic and the standard errors ignore it, if the timestamps are those
of publication rather than of the information (look-ahead), if many signals are tested and
the best one is reported, and if the final value is itself revised.
\iqlookfor{regression of revisions on prior information, and the ways the test fails.}
\end{solution}
